\chapter{연구 배경}
\section{SO-101 플랫폼과 관측 장치}
SO-101은 리더암과 팔로워암을 이용한 텔레오퍼레이션 환경을 구성할 수 있는 로봇팔이다. LeRobot의 장비 문서는 조립, 통신, 캘리브레이션 및 리더·팔로워 구성 절차를 제공한다\cite{so101}. 본 과제는 리더암으로 시연과 캘리브레이션 자세를 만들고, 팔로워암으로 실제 블록 조작을 수행하였다. 팔의 위치·방향을 제어하는 다섯 관절과 그리퍼 개폐 축을 구분하여 다룬다.

초기 모방학습 단계에서는 손목 카메라와 상단 카메라를 사용하였다. 좌표 변환과 파지 계획은 고정된 상단 관측 카메라를 기준으로 하며, 최근에는 파지 검증을 보완하기 위해 손목 카메라를 다시 장착하는 작업을 진행하고 있다(3장). 비스듬한 시점에서 생기는 작업 평면의 원근 변형은 homography로 보정한다. 평면과 높이가 다른 물체에는 잔여 투영 오차가 발생하므로, 변환의 기준 높이를 블록 윗면에 맞추었다.

\begin{figure}[htbp]
\centering
\begin{subfigure}{0.61\linewidth}
\includegraphics[width=\linewidth]{setup_early.png}
\caption{중간보고서 시점의 장비 배치}
\end{subfigure}\hfill
\begin{subfigure}{0.32\linewidth}
\includegraphics[width=\linewidth]{blocks.png}
\caption{과제 대상 블록}
\end{subfigure}
\caption{모방학습 단계의 장비 배치와 과제 대상 블록.}
\label{fig:hardware}
\end{figure}

팔로워의 관절 위치와 부하 정보는 제어와 결과 확인에 사용된다. 관절각의 피드백으로 동작 보간의 시작 자세를 정하고, 그리퍼의 위치와 부하로 빈손 닫힘과 물체를 물고 있는 상태를 구분한다. 그리퍼 위치는 장비 캘리브레이션에 따른 정규화 값이며, 부하는 힘 단위로 교정하지 않은 서보 피드백 값이다.

\section{고전 컴퓨터 비전 기반 블록 검출}
HSV 색 공간에서 OpenCV의 채널별 하한·상한 임계값으로 이진 마스크를 생성하여 블록 후보를 만든다\cite{opencv-hsv}. 빨강처럼 색상 범위의 양 끝에 걸리는 색은 두 구간의 마스크를 합쳐 처리한다. 색 임계값은 고정된 장면의 관측을 바탕으로 정한 설정이며, 조명과 배경이 달라지면 재검토가 필요하다.

색이 비슷한 배경을 배제하기 위해 윤곽선의 면적과 회전 사각형, 볼록 껍질을 함께 사용한다. 윤곽선 속성을 이용한 면적·종횡비·solidity 등의 계산은 OpenCV에서 제공하는 기본 연산으로 구성할 수 있다\cite{opencv-contour}. 윤곽선 면적을 $A$, 회전 경계 사각형의 폭과 높이를 $w,h$, 볼록 껍질의 면적을 $A_{\mathrm{hull}}$이라 하면 다음과 같이 정의한다.
\begin{equation}
 r=\frac{\max(w,h)}{\min(w,h)},\qquad
 f=\frac{A}{wh},\qquad
 s=\frac{A}{A_{\mathrm{hull}}}.
\end{equation}
여기서 $r$은 종횡비, $f$는 채움 비율, $s$는 solidity이다. 긴 테이프 조각은 종횡비가 커질 수 있고, 내부가 빈 경계 형태는 채움 비율이 낮아질 수 있다. 실제 블록도 그림자나 측면의 영향으로 이상적인 정사각형과 다르게 검출되므로, 형상 조건은 실환경 결과와 함께 조정해야 한다.

색상과 형상은 후보 생성의 기준이고, 실제 제어 대상의 결정에는 작업 범위와 배치 상태도 필요하다. 본 과제에서 사용한 색 기준점 할당, 형상 임계값과 영역 필터의 적용 순서는 3장에서 설명한다.

\section{평면 좌표 변환과 캘리브레이션}
\subsection{Homography의 적용 범위}
Homography는 두 평면 좌표 표현 사이의 사영 변환을 나타내며, 평면의 원근 변형을 보정하거나 대응점 관계를 표현하는 데 사용된다\cite{opencv-homography}. 본 시스템은 영상 픽셀 $(u,v)$와 로봇 베이스 프레임의 평면 좌표 $(x,y)$를 다음과 같이 연결한다.
\begin{equation}
 \lambda\begin{bmatrix}x\\y\\1\end{bmatrix}
 =H\begin{bmatrix}u\\v\\1\end{bmatrix},\quad
 x=\frac{h_{11}u+h_{12}v+h_{13}}{h_{31}u+h_{32}v+h_{33}},\quad
 y=\frac{h_{21}u+h_{22}v+h_{23}}{h_{31}u+h_{32}v+h_{33}}.
\end{equation}
실제 계산에서는 분모로 정규화한 좌표를 사용한다. 로봇 베이스를 원점으로 정의하여 검출 위치, 지정 영역과 거리 기반 대상 선택을 같은 좌표계에서 표현한다. 체스판은 대응점 수집과 카메라 드리프트 관찰에 사용한다.

카메라 마운트가 움직이거나 기준 높이가 바뀌면 픽셀과 로봇 좌표의 대응 관계도 달라져 재캘리브레이션이 필요하다.

\subsection{잔차와 일반화 오차}
캘리브레이션의 적합 잔차는 대응점 $\boldsymbol{p}_i$와 변환으로 예측한 위치 $\hat{\boldsymbol{p}}_i$의 차이로 계산한다. $N$개 점에 대한 제곱평균제곱근 오차는 다음과 같다.
\begin{equation}
 e_i=\|\hat{\boldsymbol{p}}_i-\boldsymbol{p}_i\|_2,\qquad
 \mathrm{RMS}=\sqrt{\frac{1}{N}\sum_{i=1}^{N} e_i^2}.
\end{equation}
같은 점으로 변환을 적합하고 평가한 잔차는 새로운 위치의 오차를 직접 나타내지 않는다. 이에 따라 한 점을 제외하고 나머지 점으로 변환을 적합한 뒤 제외한 점을 평가하는 leave-one-out(LOO) 결과를 함께 확인한다. 또한 대응점의 로봇 좌표 자체를 FK로 기록하므로, 이 평가에는 카메라뿐 아니라 자세 기록과 기구학 기준의 오차가 함께 포함될 수 있다.

\section{역기구학과 상태 기반 제어}
\subsection{목표 위치에서 관절 자세로의 변환}
역기구학은 목표 자세 $T_d$를 만족하는 관절값 $\boldsymbol{q}$를 구하는 문제이며, 수치해법에서는 초기값과 도달 가능성이 수렴 결과에 영향을 미친다\cite{modern-robotics}. 본 구현은 URDF 모델과 PlaCo의 운동학 풀이 기능을 사용한다. PlaCo는 관절 제한 등을 포함한 작업 공간의 역기구학 문제를 구성할 수 있다\cite{placo}.

SO-101에서는 위치와 모든 방향을 임의로 독립 지정할 수 없으므로, 목표 위치뿐 아니라 접근 방향과 관절 한계를 함께 검사해야 한다.

역기구학은 접근·파지·운반 등 소수의 웨이포인트에서 수행하며, 세부 보간 방식과 실행 결과 확인은 3장에서 설명한다.

\subsection{유한상태기계와 파지 확인}
유한상태기계(FSM)는 대상 선택부터 배치까지의 단계를 상태로 분리하여 실행 조건과 실패 경로를 명시한다. 파지가 확인되지 않으면 운반 단계로 넘어가지 않는 조건이 핵심이며, 세부 흐름은 3장에서 설명한다.

조작 결과를 평가할 때는 파지 확인, 운반 중 물체 유지, 배치 후 안착, 적재 후 유지를 구분해야 한다. 현재 센서 판정은 그리퍼가 물체를 물고 있는지를 확인하며, 이후의 안착과 적재 유지까지 자동으로 보장하지는 않는다. 이 구분은 실패 분석과 학습 데이터의 성공 라벨을 정의할 때에도 적용된다.

\FloatBarrier
